\section*{Conclusion générale}

Ce mémoire est parti d'un besoin concret et profondément économique : permettre à une
petite structure, Audit Énergie, de développer une activité de formation externe sans dépendre
du coût récurrent d'un formateur en direct. Cette contrainte initiale a guidé l'ensemble du
travail. À chaque étape, une solution a été essayée, ses limites ont été observées, puis une
réponse plus complète a été conçue : du cours en direct au cours audio préparé à l'avance, puis
aux essais de synthèse vocale, et enfin à une plateforme autonome combinant une pipeline de
génération de contenus et un assistant pédagogique contextualisé.

L'apport central du travail n'est pas l'utilisation isolée d'un modèle de langage ou d'un service
de synthèse vocale, mais leur \emph{orchestration} dans une chaîne fiable et observable. La
pipeline structurée, fondée sur un plan explicite, une génération section par section, des
reviews ciblées et des artefacts intermédiaires, transforme une sortie d'IA opaque en un
processus auditable. L'assistant RAG, ancré sur les documents réels de la formation, prolonge
l'expérience du cours en répondant aux questions des élèves sans tomber dans les dérives d'un
chatbot généraliste. Ensemble, ces deux briques répondent à la problématique posée :
automatiser fortement le modèle de formation tout en préservant la qualité pédagogique, la
traçabilité des contenus et la pertinence des réponses.

Ce travail comporte des limites assumées. Plusieurs indicateurs — temps réellement gagné,
coût complet, fidélité du RAG, précision de la calibration audio, satisfaction des élèves —
restent à consolider par un protocole d'évaluation plus systématique. La solution dépend par
ailleurs de services externes dont le coût et la disponibilité peuvent évoluer, et la supervision
humaine demeure nécessaire sur les points les plus sensibles. Ces limites ne remettent pas en
cause l'approche : elles tracent les pistes d'évolution, notamment la mise en place d'un banc
d'évaluation RAG, l'évaluation croisée automatique et humaine des contenus générés, et une
calibration audio mesurée plutôt qu'empirique.

Sur le plan professionnel, ce projet a permis une véritable montée en compétences :
l'analyse d'un besoin métier, la conception d'une architecture logicielle capable de gérer des
traitements longs et reprenables, l'ingénierie des systèmes d'IA (savoir quand recourir à un
prompt, à un RAG, à un enrichissement de source ou à une génération structurée) et, surtout,
la prise de recul nécessaire pour remettre en question une solution lorsqu'elle devient
difficile à contrôler. Au-delà de la plateforme elle-même, ce mémoire illustre une conviction
construite tout au long du travail : la valeur de l'intelligence artificielle générative en
entreprise ne tient pas à la puissance brute des modèles, mais à la rigueur de l'architecture
qui les encadre.
